% Generated draft: requires mathematical and source review.
\documentclass[11pt]{article}
\usepackage[margin=25mm]{geometry}
\usepackage{fontspec}
\setmainfont{Latin Modern Roman}
\usepackage{amsmath,amssymb,mathtools,unicode-math}
\setmathfont{Latin Modern Math}
\usepackage{enumitem,needspace,xurl,hyperref}
\setlength{\parindent}{0pt}
\setlength{\parskip}{5pt}
\setlength{\emergencystretch}{3em}
\newcommand{\N}{\mathbb{N}}
\newcommand{\Z}{\mathbb{Z}}
\newcommand{\Q}{\mathbb{Q}}
\newcommand{\R}{\mathbb{R}}
\newcommand{\C}{\mathbb{C}}
\newcommand{\F}{\mathbb{F}}
\newcommand{\A}{\mathbb{A}}
\newcommand{\PP}{\mathbb{P}}
\newcommand{\E}{\mathbb{E}}
\newcommand{\eps}{\varepsilon}
\newcommand{\dd}{\,\mathrm{d}}
\begin{document}
\section*{UnsolvedMath ID: 30005219}
\textbf{Source problem code:} catalogue entry
\subsection{Definitions and mathematical statement}
The supplied record does not explicitly specify the class of groups or the character theory intended. Under the usual interpretation in ordinary finite-group character theory, let \(G\) be a finite group and let \(\operatorname{Irr}(G)\) denote its irreducible complex characters. Such a character has the form \(\chi(g)=\operatorname{tr}(\rho(g))\), where \(\rho:G\to\operatorname{GL}(V)\) is an irreducible representation on a nonzero finite-dimensional complex vector space \(V\). Its field of values is
\[
\mathbb{Q}(\chi)=\mathbb{Q}\bigl(\chi(g):g\in G\bigr)\subseteq\mathbb{C},
\]
the smallest subfield of \(\mathbb{C}\) containing \(\mathbb{Q}\) and every value of \(\chi\).

The supplied question asks: what structural restrictions hold on \(\mathbb{Q}(\chi)\) as \(G\) varies and \(\chi\) ranges over \(\operatorname{Irr}(G)\)? Under the interpretation above, the scope is all finite groups and all their irreducible complex characters; no degree or prime restriction is specified.

This is a broad research question, rather than a proposed theorem with a specified conclusion. Finiteness of \(G\) and the ordinary complex-character interpretation are explicit interpretive assumptions in this draft and require confirmation from the source context. The record gives no precise meaning for the requested restrictions and no proposed necessary or sufficient conditions.
\subsection{Short English statement}
Determine what restrictions the field generated by the values of an irreducible character must satisfy. The supplied question leaves the desired restrictions unspecified; its interpretation for ordinary complex characters of finite groups requires confirmation.
\subsection{Sources}
\begin{itemize}
\item Local UnsolvedMath registry, permanent ID \(30005219\), title \emph{Fields of Values of Irreducible Characters}; exact source problem code \texttt{catalogue entry}. The record attributes the question to Oberwolfach Reports workshop discussions and names no individual proposer. It records proposed year 2022.
\item Supplied report citation: \emph{Character Theory and Categorification} (2023), Oberwolfach Reports, pp.\ 2267--2333, DOI 10.4171/owr/2022/39. Original report link: \url{https://doi.org/10.4171/owr/2022/39}. No specific page or named contributor for this question is supplied.
\item Supplied literature reference: Hung--Navarro--Tiep, fields of character values, \url{https://arxiv.org/abs/2601.15987}.
\item Supplied literature reference: Navarro--Ruhstorfer--Tiep--Vallejo, height-zero fields, \url{https://doi.org/10.1016/j.aim.2024.109589}.
\item Literature assessment dated exactly \texttt{2026-08-22}: the supplied curation labels the question open and describes it as intentionally broad. It reports classifications in important degree- and prime-restricted cases, remaining conjectures, and no general structural characterization of \(\mathbb{Q}(\chi)\). This is supplied dated curation, not an independently verified assessment. The record explicitly requests independent verification of time-sensitive status and cited claims.
\end{itemize}
\end{document}
